\chapter{Összegzés}

\section{Mit mutat be a dolgozat?}

A dolgozat egy olyan rendszert mutat be, amely valós League of Legends mérkőzésadatokból indul ki, majd ezekből játékosprofilokat, kapcsolatokat, gráfokat, klasztereket, elemzési eredményeket és végül egy szimulációs kísérletet épít. A rendszer nem csak azt nézi meg, hogy egy játékos mennyit öl vagy hányszor hal meg, hanem azt is, hogy kikkel játszik együtt, kikkel kerül szembe, milyen ismétlődő kapcsolatok rajzolódnak ki körülötte, és ezekből milyen nagyobb hálózati minták olvashatók ki.

Ha valamit le kellene szűkíteni a projekt lényegéből, az inkább az egész lánc megléte, mint bármelyik egyedi komponens. Önmagában egy Dijkstra-implementáció vagy egy React-oldal nem különösebben érdekes. De ha a nyers meccsadatoktól el lehet jutni gráfon, klasztereken és Rust-analitikán át egy reprodukálható szimulációs kísérletig -- és minden lépés ellenőrizhető és újrafuttatható --, akkor az már más jellegű eredmény.

\section{A rendszer közérthető képe}

A kérdések, amelyeket a rendszer vizsgál, nem pszichológiai állítások. Vannak-e sűrűbb csoportok a kapcsolathálóban? Vannak-e játékosok, akik sok csoport között közvetítenek? Hasonló teljesítményű játékosok kerülnek-e egymás mellé szívesebben? Másképpen viselkedik-e egy Flex Queue-ból felépített gráf, mint egy SoloQ kontroll-gráf?

Ezek mérhető, módszertanilag védett kérdések. A válaszok megmutatják, milyen hálózati szerkezetek keletkeznek adott adatgyűjtési és elemzési szabályok mellett.

\section{A fő eredmények}

A dolgozat első nagy eredménye a több-adatkészletes adatgyűjtési és feldolgozási architektúra. A \textit{flexset} és a \textit{soloq} adatkészlet külön kezelése azért fontos, mert a két queue nem ugyanazt a jelenséget méri.

A második fontos eredmény az \textit{opscore} és \textit{feedscore} metrikák használata. Ezek nem tökéletes igazságmérők, hanem értelmezhető, szerepkörérzékeny pontszámok, amelyek lehetővé teszik, hogy a teljesítménymutatók gráfanalitikai bemenetként is használhatók legyenek.

A harmadik eredmény a játékoskapcsolati gráf és annak elemzése. A \textit{flexset} adatkészleten kirajzolódó szerkezet core-periphery jellegű: vannak sűrűbb, erősebben összekapcsolt magrészek, és sok lazább, kisebb periférikus csoport.

A negyedik eredmény az asszortativitási elemzés. Az \textit{opscore} esetében pozitív kapcsolat figyelhető meg: a kapcsolt játékosok több esetben hasonlóbb teljesítményprofilt mutatnak, mint amit teljesen véletlen kapcsolatoknál várnánk.

Az ötödik eredmény a Rust-alapú algoritmikus réteg: útvonalkeresés, asszortativitás és betweenness-központiság determinisztikus és újrafuttatható módon.

\section{A Tribal NeuroSim szerepe}

A Tribal NeuroSim inkább egy transfer-kísérlet: azt vizsgálja, hogy a gráfból és a klaszterekből kinyert teljesítményprofilok átvihetők-e egy evolúciós, többügynökös szimuláció kezdeti paramétereibe.

A négy validált futtatás alapján a Flex Queue-ból származó profilok konvergensebb viselkedést mutattak: mindkét vizsgált seed mellett hasonló, Warband/Combat jellegű győztes archetípus jelent meg. A SoloQ futtatások változékonyabbak voltak. Ez azt támasztja alá, hogy a két adatkészletből származó klaszterpriorok ugyanabban a szimulációs szabályrendszerben eltérő emergáló dinamikához vezethetnek.

\section{Módszertani határok}

A signed balance -- az előjeles gráfok strukturális egyensúlya -- technikailag érdekes kísérleti irány marad, de nem fő empirikus eredmény. Ennek oka, hogy az ellenségként modellezett mérkőzéskapcsolat nem feltétlenül jelent valódi negatív társas kapcsolatot.

Ugyanezért nem szerepelnek a dolgozat fő állításai között időbeli stabilitási vagy közösségi kohézióra vonatkozó erős kijelentések. A projekt végső formája inkább a védhetőbb kérdésekre koncentrál.

\section{Végső következtetés}

A dolgozat végső állítása az, hogy kompetitív játékadatokból építhető olyan rendszer, amely egyszerre szoftvermérnöki és kutatási szempontból is értelmezhető. A PremadeGraph nem pusztán statisztikai oldal, hanem adatgyűjtési, adatmodellezési, gráfanalitikai és szimulációs keretrendszer.

Ami a rendszerből visszatekintve a legjobban működött, az a lépések közötti átjárhatóság: a nyers meccsadatokból normalizált játékosprofil lesz, abból kapcsolati gráf, a gráfból klaszterek és analitikai mutatók, a klaszterekből pedig szimulációs priorok.

\section{Túl a videojátékon: módszertani transzferálhatóság}
\label{sec:osszegzes-transfer}

A PremadeGraph módszertana nem LoL-specifikus. Amit ez a projekt épített, az egy általános célú kutatási infrastruktúra, amelynek minden komponense transzferálható.

\textbf{Gráfklaszterezés és közösségi struktúra.} A greedy modularity maximization egyaránt alkalmazható tudományos kollaborációs hálózatokra, vállalati szervezeti gráfokra, logisztikai hub-hálózatokra vagy közösségi média interakciókra.

\textbf{Asszortativitási elemzés.} Pénzügyi hálózatokban a kockázati kontagion terjedési mintáját írja le; ökológiai hálózatokban a niche-specializáció fokát méri; epidemiológiai modellekben a kontaktstruktúra homofíliáját jellemzi.

\textbf{Párhuzamos Brandes-féle betweenness centralitás.} A kritikus közvetítő csomópontok azonosítása infrastruktúra-tervezésben, járványvédelemben és ellátásilánc-elemzésben egyaránt közvetlen alkalmazást talál.

\textbf{Pathfinder algoritmusok.} A BFS, Dijkstra, kétirányú keresés és az egzakt A* navigációs rendszerek, hálózati routing optimalizáció, robotikai útvonaltervezés esetén egyaránt alkalmazhatók.

\textbf{Tribal NeuroSim evolúciós szimulációja.} Az artifact-prior alapú inicializálás és a NEAT-stílusú neuroevolúció kombinációja bármilyen domainre alkalmazható, ahol csoportszintű teljesítményprofilok kinyerhetők és összehasonlíthatók. Az öt alapartifakt (combat, resource, map\_objective, risk, team) csapatsportokban, startupok esetén és ökológiai modellekben egyaránt értelmezhető.

A Riot Games API azért volt kényelmes prototípus-forrás, mert nyilvánosan elérhető, gazdagon dokumentált és strukturált JSON adatot szolgáltat. De a módszertan nem igényel LoL-adatot. Ugyanez a pipeline futtatható bármely nyílt API-ra (GitHub kollaborációs gráf, NBA teljesítményadatbázis, PubMed szerzői hálózat), ahol hasonló felosztás elvégezhető.

Ahogy az ABBA megfogalmazta évtizedekkel ezelőtt: \textit{the winner takes it all}.
A Tribal NeuroSimben ez nem metafora, hanem a mechanika közvetlen következménye, és egyben az egész pipeline szimbóluma. Az adat, a gráf, a klaszter, az analitika és a szimuláció mind egyetlen coherens lánc részei. Aki megérti ezt a láncot, annak kezében van egy általánosan alkalmazható hálózatelemzési és evolúciós modellezési keretrendszer, amelynek a videojáték csak a legpraktikusabb és legelérhetőbb bemeneti forrása volt.
